\documentclass[11pt]{article}
\usepackage[margin=0.9in]{geometry}
\usepackage{amsmath,amssymb,amsthm,mathtools}
\usepackage[T1]{fontenc}
\usepackage[utf8]{inputenc}
\usepackage{enumitem}
\usepackage{booktabs,longtable,array,seqsplit}
\usepackage{hyperref}
\usepackage{xcolor}
\usepackage{microtype}
\hypersetup{colorlinks=true,linkcolor=blue!45!black,citecolor=blue!45!black,urlcolor=blue!45!black}
\setlength{\emergencystretch}{6em}

\newtheorem{theorem}{Theorem}[section]
\newtheorem{proposition}[theorem]{Proposition}
\newtheorem{lemma}[theorem]{Lemma}
\newtheorem{corollary}[theorem]{Corollary}
\newtheorem{claim}[theorem]{Claim}
\theoremstyle{definition}
\newtheorem{definition}[theorem]{Definition}
\newtheorem{assumption}[theorem]{Assumption}
\theoremstyle{remark}
\newtheorem{remark}[theorem]{Remark}

\newcommand{\R}{\mathbb{R}}
\newcommand{\C}{\mathbb{C}}
\newcommand{\N}{\mathbb{N}}
\newcommand{\Hsp}{\mathcal{H}}
\newcommand{\Fsp}{\mathcal{F}}
\newcommand{\Id}{\mathrm{Id}}
\newcommand{\tr}{\operatorname{tr}}
\newcommand{\diag}{\operatorname{diag}}
\newcommand{\rank}{\operatorname{rank}}
\newcommand{\spec}{\operatorname{spec}}
\newcommand{\spanop}{\operatorname{span}}
\newcommand{\dist}{\operatorname{dist}}
\newcommand{\range}{\operatorname{ran}}
\newcommand{\Ker}{\operatorname{ker}}
\newcommand{\Pinch}{\mathcal{C}}
\newcommand{\Off}{\widetilde{\mathcal{C}}}
\newcommand{\Pperp}{P^{\perp}}
\newcommand{\Qperp}{Q^{\perp}}
\newcommand{\op}{\mathrm{op}}
\newcommand{\F}{\mathrm{F}}
\newcommand{\HS}{\mathrm{HS}}
\newcommand{\ip}[2]{\left\langle #1,#2\right\rangle}
\newcommand{\norm}[1]{\left\lVert #1\right\rVert}
\newcommand{\abs}[1]{\left\lvert #1\right\rvert}
\newcommand{\code}[1]{\nolinkurl{#1}}

\newenvironment{argument}{\paragraph{Argument.}\begin{enumerate}[leftmargin=2em]}{\end{enumerate}}
\newenvironment{proofoutline}{\paragraph{Proof architecture.}\begin{enumerate}[leftmargin=2em]}{\end{enumerate}}
\newcommand{\sourceanchor}[1]{\par\smallskip\noindent\textbf{Source anchor.} #1\par\smallskip}
\newcommand{\sourcecomment}[1]{\begin{remark}#1\end{remark}}
\newenvironment{sourceguide}{\begin{description}[style=nextline,leftmargin=3.3cm,labelwidth=3.1cm]}{\end{description}}
\newenvironment{formalizationmap}{\begin{longtable}{@{}p{0.20\textwidth}p{0.31\textwidth}p{0.40\textwidth}@{}}\toprule Source item & Modern mathematical content & Repository status \\\midrule\endfirsthead\toprule Source item & Modern mathematical content & Repository status \\\midrule\endhead}{\bottomrule\end{longtable}}

\title{Motovilov (2012): graph subspaces, the relaxed tan $\Theta$ theorem, and an a priori companion}
\author{}
\date{Source-order distilled reconstruction of arXiv:1204.4441v2}
\begin{document}
\maketitle

\section*{Source map and scope}
\begin{sourceguide}
\item[Primary source] Alexander K. Motovilov,\linebreak \emph{Comment on ``The tan theta theorem with relaxed conditions''} by Y. Nakatsukasa, arXiv:1204.4441v2 (2012).
\item[Coverage] Proposition 1; Remark 2; Lemma 3; Proposition 4 and its proof; Remark 5; Proposition 6 and Remark 7.
\item[Formalization target] The no-pole content and proof interpretation of \code{tan_theta_le}; canonical graph-subspace and spectral-subspace wrappers; the distinction between a posteriori and a priori tan bounds.
\item[Norm scope] The comment deliberately works with the spectral norm and the largest principal angle. It does not reproduce Nakatsukasa's every-UI-norm theorem.
\end{sourceguide}

The comment makes two corrections of perspective rather than correcting Nakatsukasa's theorem. First, the relaxed spectral-norm theorem is a finite-dimensional instance of an earlier block-operator graph-subspace theorem. Second, the graph condition required by that theorem is automatic under the finite-dimensional spectral hypotheses. It then places the a posteriori theorem beside a sharp a priori theorem with a residual smallness threshold.

\section{Proposition 1: the block-operator tan theorem}

\sourceanchor{Proposition 1.}

Let
\[
  L=\begin{pmatrix}A_1&B^*\\B&A_2\end{pmatrix}
\]
be Hermitian on $\C^k\oplus\C^{n-k}$. Let
\[
  \mathcal A_1=\C^k\oplus\{0\},
  \qquad
  \mathcal A_2=\{0\}\oplus\C^{n-k}.
\]
Assume that for some $\alpha\le\beta$ and $\delta>0$,
\[
  \spec(A_1)
  \subseteq(-\infty,\alpha-\delta]
    \cup[\beta+\delta,\infty).
\]
Let $\mathcal L_2$ be a reducing spectral subspace of $L$ with
\[
  \spec(L|_{\mathcal L_2})\subseteq[\alpha,\beta],
\]
and let $\mathcal L_1=\mathcal L_2^\perp$ have dimension $k$. Then
\begin{equation}
  \tan\angle(\mathcal A_1,\mathcal L_1)
  \le\frac{\norm{B}_{\op}}{\delta}.
  \tag{M-2}\label{M-2}
\end{equation}
Here $\angle$ denotes the largest principal angle.

The hypothesis is exactly the relaxed interval/exterior orientation: the coordinate block $A_1$ lies outside, while the complementary exact spectral subspace lies in the middle interval.

\section{Remark 2: the graph formulation}

Earlier operator-Riccati results state the estimate using the angle between the complementary subspaces $\mathcal A_2$ and $\mathcal L_2$, provided $\mathcal L_2$ is the graph of an operator
\[
  X:\mathcal A_2\to\mathcal A_1.
\]
For a graph subspace,
\[
  \tan\angle(\mathcal A_2,\mathcal L_2)=\norm{X}_{\op}.
\]
Because complementary principal-angle lists agree, this is equivalent to \eqref{M-2} when dimensions match.

The apparent extra graph assumption is not a genuine finite-dimensional hypothesis. Lemma 3 proves it from spectral separation.

\section{Lemma 3: automatic transversality}

\sourceanchor{Lemma 3 and equation (3).}

Under Proposition 1,
\begin{equation}
  \mathcal A_2\cap\mathcal L_1=\{0\},
  \qquad
  \mathcal A_1\cap\mathcal L_2=\{0\}.
  \tag{M-transverse}
\end{equation}
Hence $\mathcal L_2$ is the graph of an operator from $\mathcal A_2$ to $\mathcal A_1$.

\begin{proofoutline}
\item Since $\dim\mathcal L_1=\dim\mathcal A_1=k$, the canonical decomposition of two finite-dimensional subspaces gives
\[
  \dim(\mathcal A_2\cap\mathcal L_1)
  =\dim(\mathcal A_1\cap\mathcal L_2).
\]
It is enough to show the second intersection is zero.
\item Suppose $0\ne y\in\mathcal A_1\cap\mathcal L_2$. Write $y=(x,0)$ with $x\ne0$.
\item Center the interval at
\[
  c=\frac{\alpha+\beta}{2}.
\]
Because $y$ occupies only the first coordinate block,
\[
  \norm{(L-cI)y}^2
  =\norm{(A_1-cI)x}^2+\norm{Bx}^2
  \ge\left(\frac{\beta-\alpha}{2}+\delta\right)^2\norm{y}^2.
\]
\item Because $y\in\mathcal L_2$ and the spectrum of $L|_{\mathcal L_2}$ lies in $[\alpha,\beta]$,
\[
  \norm{(L-cI)y}
  \le\frac{\beta-\alpha}{2}\norm{y}.
\]
\item The two inequalities contradict $\delta>0$. Therefore $y=0$, and the dimension identity gives the other intersection.
\end{proofoutline}

This is the cleanest source explanation of the ``pole-free'' tangent theorem. The denominator $\cos\Theta$ is not assumed nonzero. Spectral coercivity proves transversality, and transversality proves the graph representation.

\section{Proposition 4: Nakatsukasa's theorem as a corollary}

\sourceanchor{Proposition 4 and its proof.}

Let $A$ be Hermitian, let $Q_1$ have orthonormal columns, and set
\[
  A_1=Q_1^*AQ_1,
  \qquad
  R=AQ_1-Q_1A_1.
\]
Let $X=[X_1\ X_2]$ be a unitary eigenvector matrix for $A$. Assume
\[
  \spec(A_1)
  \subseteq(-\infty,\alpha-\delta]
      \cup[\beta+\delta,\infty),
  \qquad
  \spec(\Lambda_2)\subseteq[\alpha,\beta].
\]
Then
\begin{equation}
  \tan\angle(Q_1,X_1)
  \le\frac{\norm{R}_{\op}}{\delta}.
  \tag{M-4}
\end{equation}

Complete $Q_1$ to a unitary $Q=[Q_1\ Q_2]$. In the $Q$ coordinates,
\[
  L=Q^*AQ
  =\begin{pmatrix}
      A_1&B^*\\B&A_2
    \end{pmatrix},
  \qquad
  B=Q_2^*AQ_1=Q_2^*R.
\]
Since $R\perp\range(Q_1)$,
\[
  \norm{B}_{\op}=\norm{R}_{\op}.
\]
The subspaces
\[
  \mathcal L_1=Q^*\range(X_1),
  \qquad
  \mathcal L_2=Q^*\range(X_2)
\]
reduce $L$, and the spectrum on $\mathcal L_2$ is $\spec(\Lambda_2)$. Proposition 1 applies, and unitary invariance of principal angles transports the result back to $Q_1$ and $X_1$.

Remark 5 observes that the implication reverses: Proposition 1 is obtained from Proposition 4 by taking $Q=I$. Thus the matrix residual theorem and the block graph-subspace theorem are equivalent formulations in finite dimensions.

\section{Proposition 6: the a priori tan theorem}

\sourceanchor{Proposition 6 and Remark 7.}

Proposition 4 is a posteriori: it assumes knowledge of the exact unwanted eigenvalues $\Lambda_2$. Proposition 6 replaces this with separation of the two computable compression blocks.

Let
\[
  A_1=Q_1^*AQ_1,
  \qquad
  A_2=Q_2^*AQ_2,
  \qquad
  R=AQ_1-Q_1A_1.
\]
Assume
\[
  \spec(A_1)
  \subseteq(-\infty,a-d]\cup[b+d,\infty),
  \qquad
  \spec(A_2)\subseteq[a,b],
\]
and
\begin{equation}
  \norm{R}_{\op}<\sqrt2\,d.
  \tag{M-small}
\end{equation}
Then the eigenvalues of $A$ can be partitioned so that the exact exterior subspace $X_1$ has the same dimension as $Q_1$, the complementary spectrum remains in an explicitly enlarged interval, and
\begin{equation}
  \tan\angle(Q_1,X_1)
  \le\frac{\norm{R}_{\op}}{d}.
  \tag{M-5}
\end{equation}

The spectral enlargement is controlled by
\[
  \delta_R
  =\norm{R}_{\op}
    \tan\left(
      \frac12\arctan\frac{2\norm{R}_{\op}}{d}
    \right)
  <d.
\]
The smallness condition guarantees that the perturbed spectral cluster exists with the required cardinality and remains separated. Remark 7 warns that it cannot simply be dropped: beyond the threshold, the desired middle cluster may disappear.

This theorem has two logically distinct components:
\begin{enumerate}[leftmargin=2em]
\item spectral branch selection and cluster persistence;
\item the tangent estimate after the correct exact subspace has been identified.
\end{enumerate}
A robust formal API should keep these components separable.

\section{Relation to the coordinate-free Lean proof}

The current theorem \code{tan_theta_le} proves, for every $x$ in the trial subspace $Z$,
\[
  \delta\norm{x-P_Vx}
  \le\rho\norm{P_Vx}.
\]
Its coercivity hypothesis on the trial compression is the basis-free analogue of
\[
  \spec(A_1)
  \subseteq(-\infty,\alpha-\delta]
      \cup[\beta+\delta,\infty),
\]
while the form bounds on $V^\perp$ replace
\[
  \spec(\Lambda_2)\subseteq[\alpha,\beta].
\]
The residual bound is the normwise form of $R=AQ_1-Q_1A_1$.

The proof uses a maximal projection vector in $V^\perp$ rather than explicitly constructing the graph operator. Its key coercive contradiction is the same one as Lemma 3. The result is therefore a coordinate-free modernization of Motovilov's graph argument, not an unrelated theorem.

The conclusion directly supplies transversality:
\[
  x\in Z\cap V^\perp
  \implies
  \delta\norm{x}\le0
  \implies x=0.
\]
In finite equal dimension, this is enough to identify $Z$ as a graph over $V$ and to recover the maximal tangent bound.

\begin{formalizationmap}
Proposition 1 & Block Hermitian operator; exterior coordinate block; exact middle spectral subspace; largest-angle tan bound & Conceptual block form of \code{tan_theta_le} \\
Lemma 3 & Spectral coercivity forces both cross intersections to vanish & Pole-free content of \code{tan_theta_le}; should be available as a named transversality corollary \\
Proposition 4 & Residual/Rayleigh--Ritz form of Nakatsukasa's theorem & \code{tan_theta_le} and \code{partIII_tanTheta_vector} \\
Unitary coordinate change & $B=Q_2^*R$ and $\|B\|=\|R\|$ & Projection residual and adjoint residual lemmas in \code{TanTheta.lean} \\
Proposition 6 & A priori compression-gap theorem plus spectral branch selection & Not the same as the current endpoint; relevant to future canonical spectral wrappers and continuation theory \\
Remark 7 & Residual smallness is needed for cluster existence, not merely for numerical sharpness & Must remain explicit in any a priori wrapper \\
\end{formalizationmap}

\section{What this comment settles for the finite DK API}

\begin{enumerate}[leftmargin=2em]
\item The tan theorem can be stated without an assumed inverse cosine or assumed graph operator.
\item The correct proof order is spectral coercivity $\Rightarrow$ transversality $\Rightarrow$ graph/tangent.
\item The relaxed theorem is inherently oriented: approximate exterior versus exact interior.
\item The current source-faithful public endpoint may remain a per-vector/operator-norm theorem even though Nakatsukasa proves more for UI norms; the two sources must not be conflated.
\item Canonical spectral-subspace wrappers need a separate branch-selection theorem, especially in the a priori setting.
\end{enumerate}

\section*{Bibliography}
\begin{thebibliography}{9}
\bibitem{Mot2012}
A. K. Motovilov,
\emph{Comment on ``The tan theta theorem with relaxed conditions'' by Y. Nakatsukasa},
arXiv:1204.4441v2 (2012).

\bibitem{Nak2012}
Y. Nakatsukasa,
\emph{The tan $\theta$ theorem with relaxed conditions},
Linear Algebra Appl. 436 (2012), 1528--1534.

\bibitem{DK1970}
C. Davis and W. M. Kahan,
\emph{The Rotation of Eigenvectors by a Perturbation. III},
SIAM J. Numer. Anal. 7 (1970), 1--46.
\end{thebibliography}

\end{document}
